% =====================================================================
%  WSPÓLNA PREAMBUŁA PRACY INŻYNIERSKIEJ (PWr, WEFiM)
%  Kompilacja: LuaLaTeX + biber  (latexmk -lualatex main.tex)
%  Wymogi: A4, druk dwustronny, 12 pt, Times New Roman, tekst
%  wyjustowany, rysunki/tabele/wzory numerowane w obrębie rozdziału,
%  numer strony w górnym zewnętrznym rogu.
%  Parametry do szybkiej zmiany: sekcja „USTAWIENIA” poniżej.
% =====================================================================

% ---------------------------------------------------------------------
%  USTAWIENIA
% ---------------------------------------------------------------------
% Interlinia (setspace): 1.1 w LaTeX ≈ 1,15 w Wordzie (≈ 400 słów/str.,
% zgodnie z wspolne/budzet_stron.md). 1.0 = pojedyncza, 1.25 ≈ 1,5 w Wordzie.
\providecommand{\interlinia}{1.1}

% ---------------------------------------------------------------------
%  Język i czcionki
% ---------------------------------------------------------------------
\usepackage{etoolbox}
\usepackage{fontspec}
\usepackage{polyglossia}
\setdefaultlanguage{polish}
\setotherlanguage{english}

% Times New Roman (Windows/Overleaf), w przeciwnym razie metryczny klon
% TeX Gyre Termes. Matematyka: TeX Gyre Termes Math (spójna z Times).
\usepackage{amsmath,amssymb}
\usepackage{unicode-math}
\IfFontExistsTF{Times New Roman}
  {\setmainfont{Times New Roman}}
  {\setmainfont{TeX Gyre Termes}}
\setmathfont{TeX Gyre Termes Math}
\setsansfont{TeX Gyre Heros}
\setmonofont{TeX Gyre Cursor}[Scale=MatchLowercase]
\usepackage{microtype}
\usepackage{indentfirst}   % wcięcie także pierwszego akapitu (typografia polska)

% ---------------------------------------------------------------------
%  Układ strony
% ---------------------------------------------------------------------
% Marginesy wg Golden Standard: 25 mm, od strony grzbietu 30 mm.
\usepackage[a4paper, inner=30mm, outer=25mm, top=25mm, bottom=25mm,
            headheight=15pt, headsep=6mm, footskip=12mm]{geometry}
\usepackage{setspace}
\setstretch{\interlinia}
\frenchspacing
\clubpenalty=10000 \widowpenalty=10000 \displaywidowpenalty=10000

% Numer strony w górnym zewnętrznym rogu (parzyste: lewy, nieparzyste:
% prawy), bez żywej paginy i bez stopki. Tak samo na stronach
% otwierających rozdziały.
\usepackage{fancyhdr}
\fancypagestyle{praca}{%
  \fancyhf{}%
  \fancyhead[LE,RO]{\thepage}%
  \renewcommand{\headrulewidth}{0pt}%
  \renewcommand{\footrulewidth}{0pt}}
\fancypagestyle{plain}{%
  \fancyhf{}%
  \fancyhead[LE,RO]{\thepage}%
  \renewcommand{\headrulewidth}{0pt}}
\pagestyle{praca}
% mwrep używa stylów „opening”/„closing”/„blank” na początku i końcu
% rozdziału oraz na pustych stronach — wszędzie ten sam numer strony.
\makeatletter
\let\ps@opening\ps@praca
\let\ps@closing\ps@praca
\let\ps@blank\ps@praca
\makeatother

% Spis treści do poziomu podrozdziału (x.y.z)
\setcounter{tocdepth}{2}
\setcounter{secnumdepth}{3}

% ---------------------------------------------------------------------
%  Matematyka, jednostki, chemia
% ---------------------------------------------------------------------
% Numeracja wzorów (r.n) w obrębie rozdziału, po prawej
\numberwithin{equation}{chapter}

\usepackage{siunitx}
% Jedno ustawienie dla całej pracy (słownik, p. 3.4): przecinek
% dziesiętny, mnożenie „·”, zakresy z półpauzą i jednostką raz na końcu.
\sisetup{
  output-decimal-marker = {,},
  exponent-product      = \cdot,
  group-separator       = {\,},
  group-minimum-digits  = 4,
  range-phrase          = {--},
  range-units           = single,
  list-pair-separator   = {~i~},
  list-final-separator  = {~i~},
  list-separator        = {,~},
  per-mode              = symbol,
  inter-unit-product    = \ensuremath{{}\cdot{}},
  separate-uncertainty  = true,
}
% Wsteczna zgodność: istniejące rozdziały używają \SI, \SIrange, \si, \num
% (składnia siunitx v2) — działa bez zmian. Nowy tekst: \qty, \qtyrange.

\usepackage[version=4]{mhchem}    % \ce{SiO2}, \ce{Si3N4} (opcjonalnie)

% ---------------------------------------------------------------------
%  Rysunki, tabele, wykresy, schematy
% ---------------------------------------------------------------------
\usepackage{graphicx}
\usepackage{float}
\usepackage{booktabs}
\usepackage{multirow}
\usepackage{array}
\usepackage{tabularx}
\usepackage{enumitem}
\setlist{nosep, leftmargin=*, labelindent=\parindent}
\setlist[itemize]{label=\textbullet}  % kropki zamiast półpauz (prośba autora 09.11.2026)

\usepackage{xcolor}
\usepackage{tikz}
\usetikzlibrary{arrows.meta, positioning, calc, decorations.pathmorphing}
\usepackage[european, siunitx]{circuitikz}
\usepackage{pgfplots}
\pgfplotsset{compat=1.17,
  /pgf/number format/use comma,
  /pgf/number format/1000 sep={}}

% Fragmenty-rysunki zapisane jako samodzielne dokumenty (z \documentclass
% i preambułą, np. uklad-pomiarowy/uklad_pomiarowy_body.tex) wstawia się
% makrem \inputBezPreambuly{plik} z szablon/makra.tex.

% Podpisy: „Rys. 3.1.” pod rysunkiem, „Tab. 3.1.” nad tabelą
\usepackage{caption}
\usepackage{subcaption}
\captionsetup{labelsep=period, font=small, labelfont=bf, justification=centering,
              singlelinecheck=true, skip=6pt}
\captionsetup[table]{position=top}
\captionsetup[sub]{labelformat=parens, labelsep=space, labelfont=normalfont}
\gappto\captionspolish{%
  \renewcommand{\figurename}{Rys.}%
  \renewcommand{\tablename}{Tab.}%
  \renewcommand{\bibname}{Literatura}%
  \renewcommand{\appendixname}{Załącznik}}
% numeracja rysunków i tabel w obrębie rozdziału (np. Rys. 3.1)
\counterwithin{figure}{chapter}
\counterwithin{table}{chapter}

% Ścieżki szukania grafik (względem wspolne/). Rozdziały mogą dopisywać
% własne przez \graphicspath lub makro \figdir.
\graphicspath{{rysunki/}{../}}

% ---------------------------------------------------------------------
%  Bibliografia: biblatex + biber, styl IEEE (wspolne/README.md, sekcja 3)
% ---------------------------------------------------------------------
\usepackage{csquotes}
\usepackage[backend=biber, style=ieee, sorting=none, maxnames=3, minnames=1,
            doi=true, url=false, isbn=false]{biblatex}
\addbibresource{bibliografia.bib}

% ---------------------------------------------------------------------
%  Hiperłącza (na końcu preambuły)
% ---------------------------------------------------------------------
\usepackage[hidelinks, unicode, bookmarksnumbered, pdfusetitle]{hyperref}
\usepackage{bookmark}
% tylda i twarda spacja w zakładkach PDF (tytuły z „i~”, „w~”)
\pdfstringdefDisableCommands{\def~{ }\def\nobreakspace{ }\def\TODO#1{}}
